\section{Global uniqueness of the matching parameter}\label{sec:uniqueness}
\subsection{Rational coordinates and reversibility}
The phase variables in this section are \(\zeta,\eta\). They are independent
and must not be confused with the amplitude ratio \(u\) in \(F_s\).
Write \(B_0(x)=b(x)^2=(1-x^2)/4\). On a valid complete chain set
\begin{equation}\label{eq:Phi}
 \Phi_s(x,y,u)=(x,y,q),\qquad
 q=s-d(x,y)-b(x)/u,\qquad \bar q=s-d(x,y)-q.
\end{equation}
Here \(q=b(y)v>0\), \(\bar q=b(x)/u>0\), and the inverse is
\(\Phi_s^{-1}(x,y,q)=(x,y,b(x)/\bar q)\).
Thus the relevant rational domain is \(|x|,|y|<1\), \(q,\bar q>0\).
We do not assert that every state in the larger domain \(\mathcal D\)
has positive outgoing \(q\).

\begin{lemma}[Exact conjugacy, reverser and contact identities]\label{lem:rational}
In these coordinates the propagation map and its inverse are
\begin{align}
 \widehat F_s(x,y,q)&=(y,z,Q),\qquad
 z=\frac{y}{2q}-\frac{(x-1/2)B_0(y)}{q^2}-\frac12,\label{eq:Fhat}\\
 Q&=s-d(y,z)-\frac{B_0(y)}q,\notag\\
 \widehat F_s^{-1}(x,y,q)&=(w_-,x,B_0(x)/\bar q),\notag\\
 w_-&=\frac{x}{2\bar q}
       -\frac{(y+1/2)B_0(x)}{\bar q^2}+\frac12.\notag
\end{align}
The involution
\begin{equation}\label{eq:reverser}
 R_s(x,y,q)=(-y,-x,\bar q)
\end{equation}
satisfies \(R_s\widehat F_sR_s=\widehat F_s^{-1}\), exchanges \(q,\bar q\),
and preserves decreasing label orientation under reversal of chains.
For fixed \(s\), the contact form \(\omega=dq+(x-1/2)\,dy\) satisfies
\begin{equation}\label{eq:contact}
 \widehat F_s^*\omega=\frac{B_0(y)}{q^2}\omega,\qquad
 R_s^*\omega=-\omega,\qquad
 \det D\widehat F_s=\left(\frac{B_0(y)}{q^2}\right)^2.
\end{equation}
\end{lemma}
\begin{proof}
Substitute \(v=q/b(y)\) in \eqref{eq:F}. At its image the complementary
coordinate is \(B_0(y)/q\); solving the resulting equations gives both
inverse identities. The equality \(d(-y,-x)=d(x,y)\) then proves
\eqref{eq:reverser}. For the contact calculation,
\[
 dQ+(y-1/2)\,dz
 =\left(-(z+1/2)+\frac{y}{2q}\right)dy+\frac{B_0(y)}{q^2}\,dq;
\]
the formula for \(z\) makes its \(dy\) coefficient
\((x-1/2)B_0(y)/q^2\). Differentiation also gives the stated determinant.
Appendix~\ref{app:symbolic} specifies the exact polynomial-cancellation
checks, including comparison with the implemented map and Jacobian.
\end{proof}

Reversal takes the intrinsic positive unstable tail to the intrinsic
negative stable tail. This describes the second tail as a set; it does not
assume that a given orbit meets its own reflection.

\subsection{An exactly defined analytic tail}
Put
\begin{align}
 G(t)&=\frac{(2t-1)(t+1)}{2(2t+1)},&
 p_t&=(t,t,G(t)),\label{eq:G}\\
 A(t)&=(t+1)(2t-1)^2,&
 C(t)&=(1-t)(2t+1)^2,&
 D(t)&=4t^2(4t^2-1),\notag\\
 \mu(t)&=\frac{D(t)+\sqrt{D(t)^2-4A(t)C(t)}}{2A(t)},&
 \rho(t)&=\frac{C(t)}{A(t)}.\label{eq:mu}
\end{align}
At \(p_t\), the characteristic polynomial is
\[
 (\xi-\rho)\bigl(\xi^2-(D/A)\xi+\rho\bigr).
\]
Thus its eigenvalues are \(\mu,\rho,\rho/\mu\), with
\(\mu>1\) and \(0<\rho,\rho/\mu<1\) on \(T\). Direct cancellation gives
\(G(t)/b(t)=r(t)\), so this is the same positive outer fixed point.

Define \(P(t,\zeta)\) as the analytic invariant germ specified by
\begin{equation}\label{eq:germ}
\begin{gathered}
 P(t,0)=p_t,\qquad P_\zeta(t,0)=(-1/\mu,-1,t-1/2),\\
 P(t,\mu\zeta)=\widehat F_{\sigma(t)}(P(t,\zeta)).
\end{gathered}
\end{equation}
The normalization fixes the phase scale; it is part of the exact definition
of the scalar residual, not a fitted numerical convention.

\begin{certificate}[Analytic-germ enclosure]\label{cert:local}
Let \(T^\Box\) be the closed outward interval specified in
Appendix~\ref{app:intervalspec}, which contains \([a(L),a(U)]\).
On a complex parameter rectangle of half-width \(10^{-5}\) about its
midpoint and on the phase disk of radius \(r_0=0.01\), the degree-36
coefficient recursion for \eqref{eq:germ} has a unique analytic correction.
The degree-48 recursion encloses the same germ. With the sup norm on its
three coordinates, the corrections satisfy respectively
\[
 \|P-P_{36}\|<3.414\times10^{-47},\qquad
 \|P-P_{48}\|<8.334\times10^{-63}.
\]
The parameter margin exceeds \(9.996\times10^{-6}\); division denominators
are separated from zero. The contraction factors are less than
\(6.891\times10^{-27}\) and \(4.409\times10^{-36}\), respectively.
The exact inequalities, coefficient boxes and margins are the outputs of
the local certificates specified in Appendix~\ref{app:intervalspec}.
\end{certificate}
\begin{proof}[Analytic implication of the certificate]
Write \(P_N=\sum_{n=0}^Nc_n(t)\zeta^n\). Once \(c_0,c_1\) are fixed by
\eqref{eq:germ}, set the trial coefficient \(c_n\) to zero and compute
\(b_n=[\zeta^n]\widehat F_{\sigma(t)}(P_{<n})\). Then
\[
 (\mu^n\Id-J)c_n=b_n,\qquad
 J=D\widehat F_{\sigma(t)}(p_t),\qquad n\ge2.
\]
The spectral separation makes these equations unique; interval Cramer's
rule additionally rejects a determinant enclosure containing zero.
This recursion is regenerated over the complex rectangle, with the
positive-real-part square-root branch explicitly certified.

Let \(R_0=0.2\) and
\(D_N(\zeta)=\widehat F_{\sigma(t)}(P_N(\zeta/\mu))-P_N(\zeta)\).
The first \(N+1\) coefficients of \(D_N\) vanish. A coefficient majorant
bounds its norm by \(M\) on \(|\zeta|\le R_0\); hence on the smaller disk
\[
 \|D_N\|\le
 \delta_N:=M\frac{(r_0/R_0)^{N+1}}{1-r_0/R_0}.
\]
For analytic corrections vanishing to order \(N+1\), Schwarz's estimate
gives \(\|e(\zeta/\mu)\|\le\|e\|/m^{N+1}\), where the certified lower
bound is \(m>5.8365\). An enclosed Jacobian row-sum norm \(L_F\) gives
the contraction factor \(\kappa=L_F/m^{N+1}\) for the correction map
\[
 (\mathcal T_N e)(t,\zeta)=
 \widehat F_{\sigma(t)}
 \bigl(P_N(t,\zeta/\mu(t))+e(t,\zeta/\mu(t))\bigr)-P_N(t,\zeta).
\]
It preserves the required vanishing order by the coefficient equations.
The finite gates are
\(\kappa<1/2\) and \(\delta_N+\kappa\epsilon_N\le\epsilon_N\).
They make the correction map a self-map and contraction on the ball of
radius \(\epsilon_N\). Banach's theorem yields the exact analytic germ;
conjugation and uniqueness give reality on real arguments.

Cauchy estimates give parameter and phase derivative errors bounded by
\(\epsilon_N/\Delta_t\) and \(\epsilon_N/(r_0-|\zeta|)\). These errors
are added before chart propagation. All subsequent interval statements
therefore concern \(P\), not only its polynomial truncation.
The germ is injective since \(y_\zeta(t,0)=-1\), and backward iteration
divides its phase by \(\mu\). Its tangent and backward convergence, together
with Lemma~\ref{lem:tails}, identify it with the same intrinsic unstable
manifold. Its negative phase side has \(x-y<0\) near zero and cannot
represent a decreasing outer tail. The positive side does.
Finally the two degrees have identical uniquely determined Taylor
coefficients at every order; their analytic germs and valid continuations
coincide.
\end{proof}

The contact relation has a useful exact consequence. For
\(f(\zeta)=\omega(P_\zeta(\zeta))\), invariance gives
\[
 f(\zeta)=\frac{B_0(y(\zeta/\mu))}{q(\zeta/\mu)^2\mu}
                f(\zeta/\mu).
\]
The multiplier at zero is \(\rho/\mu<1\). Iteration and boundedness imply
\(f=0\) on a small germ, and analytic continuation through valid iterates
gives
\begin{equation}\label{eq:contacttangent}
 q_\zeta=(1/2-x)y_\zeta.
\end{equation}
This identity is proved analytically, rather than inferred from a small
numerical defect.

\subsection{Completeness of the two-phase domain}
Define the propagated chart
\begin{equation}\label{eq:W}
 W(t,\zeta)=\widehat F_{\sigma(t)}^{\,5}P(t,\zeta)
             =(x(t,\zeta),y(t,\zeta),q(t,\zeta)),
 \quad
 \zeta_a=0.000078125,\quad \zeta_b=0.0025.
\end{equation}
The phase interval is generated from \(r_0/128,r_0/4\), and the iterate
count is the first nonnegative count for which the right endpoint has
strictly negative \(x\). The frozen computation gives five.

\begin{certificate}[Complete chart geometry]\label{cert:geometry}
For all \(t\in T^\Box\), \(\zeta\in[0,\zeta_b]\), and all iterates
\(0,\ldots,5\), the analytic chart obeys \(|x|,|y|<1\),
\(q,\bar q>0\), and \(x_\zeta,y_\zeta<0\). At the two envelope endpoints,
\[
 y(t,\zeta_a)>0,\quad x(t,\zeta_b)<0,\quad
 x(t,\zeta_a)+y(t,\zeta_b)>0.
\]
The geometry certificate covers the closed phase interval with 14 leaves.
Every leaf contains enclosures for each intermediate iterate; its parameter
domain is the whole \(T^\Box\).
\end{certificate}
The certificate is a finite interval statement about the corrected chart
of Certificate~\ref{cert:local}. Its verification checks the signs of
interval endpoints and exact joins in the subdivision tree
(Appendix~\ref{app:intervalspec}).

\begin{lemma}[All intrinsic matches enter the phase envelope]\label{lem:coverage}
Every match in \(\mathcal M_s\), after \(\Phi_s\), satisfies
\begin{equation}\label{eq:independentmatch}
 W(t,\zeta)=R_{\sigma(t)}W(t,\eta),\qquad
 t=a(s),\quad (\zeta,\eta)\in[\zeta_a,\zeta_b]^2.
\end{equation}
The phases are independent. This closed product covers the canonical
exact-zero crossing representatives.
\end{lemma}
\begin{proof}
The functional equation gives
\[
 \widehat F_{\sigma(t)}W(t,\zeta/\mu)=W(t,\zeta),\qquad
 x(t,\zeta)=y(t,\zeta/\mu).
\]
Since \(y\) decreases strictly and \(\mu>1\), \(x>y\) for every positive
phase. The same identity applies to each earlier iterate. Certificate
\ref{cert:geometry} thus proves validity and decreasing orientation on
the entire positive propagated chart.

There are unique zeros \(\zeta_y,\zeta_x\in(\zeta_a,\zeta_b)\) of
\(y,x\), and scaling gives \(\zeta_x=\mu\zeta_y\). The precise source
crossing domain is \((\zeta_y,\zeta_x]\). Local capture identifies any
intrinsic decreasing outer tail with a positive phase of \(P\), which can
be made arbitrarily small by moving backward. Successive steps multiply
the phase by \(\mu\). A first step past \(\zeta_b\) cannot occur before
crossing: its preceding phase would exceed
\(\zeta_b/\mu>\zeta_y\), where \(y\) is already negative. Thus the first
crossing lies in \((\zeta_y,\mu\zeta_y]\), with every preceding iterate
certified. Strict decrease precludes a second crossing. This also rules
out an excursion outside the certified chart before the representative.

Apply the same reasoning independently to the reversed stable tail.
Reversal exchanges the positive \(q,\bar q\) coordinates and preserves
validity and decreasing orientation. If the canonical representative
has \(x=0>y\), its reflection has positive first label and second label
zero. Accordingly the reflected target domain is
\([\zeta_y,\zeta_x)\), whereas the source domain is
\((\zeta_y,\zeta_x]\). Both lie in the closed product used in
\eqref{eq:independentmatch}. No exact-zero state is lost.
\end{proof}

\subsection{Projected inversion and exclusion of asymmetric full matches}
The first label equation in \eqref{eq:independentmatch} is
\begin{equation}\label{eq:project}
 x(t,\eta)+y(t,\zeta)=0.
\end{equation}
For a possible source crossing, \(x(t,\zeta)\ge0\) and
\(y(t,\zeta)\le0\). The endpoint and monotonicity gates give
\[
 x(t,\zeta_a)+y(t,\zeta)>0,\qquad
 x(t,\zeta_b)+y(t,\zeta)<0,\qquad x_\eta<0.
\]
The intermediate value theorem and strict monotonicity therefore give one
and only one target phase \(\eta=\eta(t,\zeta)\) in the entire envelope.
In particular, eliminating this phase does not assume a branch near a
candidate point. The inverse is analytic wherever used, and
\begin{equation}\label{eq:inversepartial}
 \eta_\zeta=-\frac{y_\zeta(t,\zeta)}{x_\eta(t,\eta)},\qquad
 \eta_t=-\frac{y_t(t,\zeta)+x_t(t,\eta)}{x_\eta(t,\eta)}.
\end{equation}
On this projected graph, define the two remaining residuals
\begin{align}
 \HH(t,\zeta)&=x(t,\zeta)+y(t,\eta(t,\zeta)),\label{eq:HK}\\
 \KK(t,\zeta)&=q(t,\zeta)+q(t,\eta(t,\zeta))
                -\sigma(t)+d(x(t,\zeta),y(t,\zeta)).\notag
\end{align}
Projection together with \(\HH=\KK=0\) is equivalent to the complete
three-coordinate equation: if \(\HH=0\), both label pairs are reflected
and \(d(x(t,\eta),y(t,\eta))=d(x(t,\zeta),y(t,\zeta))\); then \(\KK=0\)
is exactly the complementary-\(q\) equation. Neither label equation alone
checks the amplitude ratio.

The contact identity \eqref{eq:contacttangent} and the projected equation,
without assuming \(\HH=0\), give
\begin{equation}\label{eq:HKderivatives}
 \HH_\zeta=x_\zeta(t,\zeta)
   -\frac{y_\eta(t,\eta)y_\zeta(t,\zeta)}{x_\eta(t,\eta)},\qquad
 \KK_\zeta=(y(t,\zeta)+1/2)\HH_\zeta .
\end{equation}
These exact identities control interval overestimation on projected cells.
Parameter derivatives instead use the full chain rule
\eqref{eq:inversepartial}; the fixed-parameter contact relation is not
applied to differentiation in \(t\).

\begin{certificate}[Exhaustive full-match exclusions]\label{cert:cover}
Let \(C_0=[0.0005,0.0007]\). Over the entire parameter box \(T^\Box\)
and all independent target phases in \([\zeta_a,\zeta_b]\), a closed
subdivision of the source envelope has 179 leaves. Each leaf passes one
of these predicates:
\begin{enumerate}
\item \(x<0\) or \(y>0\) throughout the cell (7 leaves);
\item the enclosure of \(\HH\) excludes zero (148 leaves);
\item the enclosure of \(\KK\) excludes zero (22 leaves);
\item the source cell lies in \(C_0\) and \(\HH_\zeta>0\)
throughout its projected family (2 leaves).
\end{enumerate}
A separate closed cover of \emph{all} of \(C_0\), with 554 leaves,
proves \(\HH_\zeta>0\) there. On \(C_0\), \(g=x+y\) has
\(g_\zeta<0\), with \(g(t,0.0005)>0\) and \(g(t,0.0007)<0\).
There are no unresolved leaves in either tree.
\end{certificate}

The finite input to each predicate is the analytic chart enclosure, not
sampled points. The inverse enclosure starts from the entire target
envelope and retains every solution under interval Newton intersection:
\[
 V_{\mathrm{new}}
 =V\cap\left(m-\frac{x(t,m)+y(t,\zeta)}{[x_\eta](t,V)}\right).
\]
The first derivative interval covers the full envelope. A floating-point
scout selects only a center, which is checked within the envelope; it does
not supply an exclusion or a target-phase restriction. Each retained
family has certified strict endpoint signs for the inverse. Subsequent
derivative intersections use already justified bounds. Tree acceptance
requires the children to have precisely the parent's endpoints and a
common join, every child to be reached exactly once, and every terminal
predicate to be resolved. These specifications are detailed in
Appendix~\ref{app:intervalspec}.

\begin{proposition}[Diagonal phases are a conclusion]\label{prop:diagonal}
For each \(t\in T\) there is a unique analytic phase
\(\phi(t)\in C_0\) with \(x(t,\phi)+y(t,\phi)=0\).
Every intrinsic full match has
\(\zeta=\eta=\phi(t)\).
\end{proposition}
\begin{proof}
The signs and nonzero phase derivative of \(g\) give its unique root in
\(C_0\), and the analytic implicit function theorem gives \(\phi\).
At that root \(\eta=\phi\) solves \eqref{eq:project}; uniqueness of
projected inversion identifies it as the target phase. Hence
\(\HH(t,\phi)=0\). The separate full central cover proves strict
monotonicity of \(\HH\) throughout \(C_0\), making \(\phi\) its only zero
there. The full cover excludes every noncentral full match by one of
the first three predicates. Lemma~\ref{lem:coverage} makes that cover
exhaustive for the intrinsic matching set, including both zero boundaries.
Thus every full match has the stated equal phases.
\end{proof}
Asymmetric solutions of the label equations are not discarded by
symmetry. The \(\KK\)-exclusion cells rigorously exclude the remaining
off-central label branches by their third-coordinate mismatch.

\subsection{Scalar equivalence and full-interval monotonicity}
We can now define
\begin{equation}\label{eq:residual}
 \Res(t)=2q(t,\phi(t))-\sigma(t)
           +d(x(t,\phi(t)),y(t,\phi(t))).
\end{equation}
Equations \eqref{eq:germ}, \eqref{eq:W}, the unique zero defining
\(\phi\), and \eqref{eq:residual} are an exact structural definition.
Truncated polynomials serve only to enclose this analytic function.

\begin{proposition}[Scalar/full-match equivalence]\label{prop:scalar}
For every \(t\in T\),
\[
 \Res(t)=0\quad\Longleftrightarrow\quad
 \mathcal M_{\sigma(t)}\ne\varnothing .
\]
\end{proposition}
\begin{proof}
Every full match is diagonal by Proposition~\ref{prop:diagonal}; its
third-coordinate equation is precisely \(\Res=0\).
Conversely, \(\Res(t)=0\) and \(g(t,\phi)=0\) give
\(x=-y\) and \(2q=\sigma-d(x,y)\). Thus \(W(t,\phi)=R_{\sigma(t)}W(t,\phi)\)
in all three coordinates. The positive-phase source tail and each of its
five propagated steps are valid and decreasing. Reflecting this chain
gives a valid decreasing forward tail to the negative fixed point.
Since \(x>y\) and \(x=-y\), the joined state has \(x>0>y\).
Applying \(\Phi_{\sigma(t)}^{-1}\) recovers every amplitude ratio and
places the join in the canonical \(\Sigma\), \(\mathcal U_s\) and
\(\mathcal V_s\). This is a genuine member of \(\mathcal M_s\).
\end{proof}

Implicit differentiation gives
\(\phi'=-(x_t+y_t)/(x_\zeta+y_\zeta)\).
Using \eqref{eq:contacttangent} in the full derivative of
\eqref{eq:residual} cancels the phase terms and yields
\begin{equation}\label{eq:residualderivative}
 \Res'(t)=2q_t-\sigma'(t)+(2x-1)y_t
 \quad\hbox{at }(t,\phi(t)).
\end{equation}
The partial derivatives on the right hold the phase fixed.

\begin{certificate}[Global derivative and endpoint signs]\label{cert:scalar}
On the entire closed box \(T^\Box\), the corrected analytic chart and
the interval enclosure of the whole phase family give
\[
 3.4800<\Res'(t)<3.5707.
\]
The analytic extension through \(a(L)\) satisfies
\[
 \Res(a(L))\in[-8.905044,-8.905043]\times10^{-16},\qquad
 \Res(a(U))\in[2.4652543,2.4652544]\times10^{-8}.
\]
These displayed outward bounds contain the full exact decimal intervals
in \pkg{scalar.json}; they are checked as rational comparisons.
\end{certificate}
This certificate encloses \(\phi(T^\Box)\) as an interval family and
evaluates \eqref{eq:residualderivative} on the full parameter interval.
It is not a set of derivative samples. Appendix~\ref{app:intervalspec}
gives the full derivative receipt and its finite acceptance rule.

\begin{theorem}[Global matching-parameter uniqueness]\label{thm:O6}
The projection
\[
 \pi_s(\mathcal M)=\{s\in I:\mathcal M_s\ne\varnothing\}
\]
is a singleton \(\{s_*\}\).
\end{theorem}
\begin{proof}
The strict derivative makes \(\Res\) increasing on \(T^\Box\), so it
has at most one zero there. Proposition~\ref{prop:scalar} and the
strict increase of \(\sigma\) then show that any two matching parameters
in \(I\) are equal. The strict endpoint signs and continuity give a zero
strictly between \(a(L)\) and \(a(U)\). The analytic extension through
the excluded left endpoint justifies this use of the intermediate value
theorem. Scalar equivalence supplies an intrinsic full match at that
zero. This existence calculation is independent of using
Proposition~\ref{prop:existmatch} as an existence premise, and uses no
fine root enclosure as an input.
\end{proof}

